\documentclass[lettersize,journal]{IEEEtran}
\usepackage{amsmath,amsfonts}
\usepackage{algorithmic}
\usepackage{algorithm}
\usepackage{array}
\usepackage[caption=false,font=normalsize,labelfont=sf,textfont=sf]{subfig}
\usepackage{textcomp}
\usepackage{stfloats}
\usepackage{url}
\usepackage{verbatim}
\usepackage{graphicx}
\usepackage{cite}
\usepackage{booktabs}
\usepackage{multirow}

\begin{document}

\title{Bridging Real-World Routing and Multi-Agent Coordination: A Hybrid Neural Combinatorial Optimization Framework}

\author{Your Name
\thanks{Department of Computer Science, Your University. Email: your.email@university.edu}}

\maketitle

\begin{abstract}
Neural Combinatorial Optimization (NCO) has shown remarkable progress in solving vehicle routing problems, yet two critical gaps persist: (1) most methods rely on simplified Euclidean distances that fail to capture real-world road network complexities, and (2) multi-agent coordination mechanisms remain underexplored despite their practical importance in fleet management. We present a hybrid framework that synergistically combines real-world routing awareness with parallel multi-agent decision-making for the Heterogeneous Capacitated Vehicle Routing Problem (HCVRP). Our approach integrates RRNCO's attention-free encoder with contextual gating and neural adaptive biases to process asymmetric distance matrices from real road networks, while leveraging PARCO's communication layers and conflict-aware parallel decoding for efficient multi-agent coordination. Experiments on 100 real-world cities demonstrate that our hybrid model achieves superior generalization on out-of-distribution test cities while maintaining computational efficiency through parallel action selection. The framework successfully bridges the gap between synthetic training environments and real-world deployment scenarios, advancing the state-of-the-art in practical neural combinatorial optimization.
\end{abstract}

\begin{IEEEkeywords}
Neural Combinatorial Optimization, Vehicle Routing Problem, Multi-Agent Systems, Real-World Routing, Attention-Free Transformers, Reinforcement Learning
\end{IEEEkeywords}

\section{Introduction}

\IEEEPARstart{V}{ehicle} routing problems (VRPs) represent a cornerstone of combinatorial optimization with profound implications for logistics, transportation, and supply chain management. Recent advances in Neural Combinatorial Optimization (NCO) have demonstrated that deep learning models can learn effective heuristics for VRPs through reinforcement learning, often matching or exceeding traditional optimization methods in solution quality while offering superior computational efficiency at inference time \cite{kool2018attention,kwon2020pomo,kim2022symnco}.

However, two fundamental limitations hinder the deployment of NCO methods in real-world scenarios. First, the overwhelming majority of existing approaches operate on simplified Euclidean distance metrics, assuming direct line-of-sight travel between locations. This abstraction fails to capture the complexities of actual road networks, including asymmetric travel times due to one-way streets, traffic patterns, elevation changes, and infrastructure constraints. Second, while real-world logistics operations typically involve coordinating multiple heterogeneous vehicles with varying capacities and speeds, most NCO research focuses on single-agent sequential decision-making, leading to inefficient autoregressive decoding that scales poorly with fleet size.

Recent works have begun to address these challenges independently. The Real Routing NCO (RRNCO) framework \cite{son2026rrnco} introduced attention-free architectures with neural adaptive biases to process real-world distance and duration matrices from OpenStreetMap data, demonstrating improved generalization to unseen cities. Concurrently, the Parallel AutoRegressive Combinatorial Optimization (PARCO) framework \cite{berto2025parco} proposed communication-enhanced parallel decoding mechanisms that enable simultaneous action selection across multiple agents, significantly reducing inference time while maintaining solution quality through priority-based conflict resolution.

\subsection{Our Contribution}

We present the first framework that synergistically combines real-world routing awareness with efficient multi-agent coordination for heterogeneous vehicle routing. Our key contributions are:

\begin{itemize}
\item \textbf{Hybrid Architecture Design}: We propose a novel integration strategy that replaces standard transformer encoders with RRNCO's attention-free network while preserving PARCO's multi-agent communication and parallel decoding mechanisms. This design leverages the complementary strengths of both frameworks—structural feature modeling from real road networks and efficient multi-agent coordination.

\item \textbf{Dual-Stream Fusion Mechanism}: We introduce a contextual gating-based fusion layer that combines coordinate-based and distance-based embeddings from RRNCO's dual-stream encoder into a unified representation compatible with PARCO's multi-agent decoder, enabling seamless information flow across architectural components.

\item \textbf{Real-World Multi-Agent Dataset}: We extend the RRNCO dataset to the multi-agent setting by generating heterogeneous fleet configurations (varying capacities and speeds) for 100 cities worldwide, providing a benchmark for evaluating NCO methods on realistic multi-agent routing scenarios.

\item \textbf{Comprehensive Evaluation}: Through extensive experiments on both in-distribution and out-of-distribution test cities, we demonstrate that our hybrid approach achieves superior generalization compared to baselines trained on synthetic Euclidean data, while maintaining the computational efficiency of parallel decoding.
\end{itemize}

The remainder of this paper is organized as follows: Section II reviews related work in neural combinatorial optimization, multi-agent routing, and real-world distance modeling. Section III details our hybrid architecture and training methodology. Section IV presents experimental results and ablation studies. Section V concludes with discussions and future directions.

\section{Related Work}

\subsection{Neural Combinatorial Optimization}

The application of deep learning to combinatorial optimization has evolved rapidly since the seminal work of Vinyals et al. \cite{vinyals2015pointer}, who introduced the Pointer Network architecture. Kool et al. \cite{kool2018attention} advanced the field with the Attention Model, demonstrating that transformer-based encoders combined with attention-based decoders could effectively solve TSP and VRP variants through REINFORCE training. Subsequent works have focused on improving sample efficiency through symmetry exploitation (SymNCO \cite{kim2022symnco}), enhancing exploration via POMO \cite{kwon2020pomo}, and incorporating problem-specific inductive biases.

However, these methods predominantly assume Euclidean distance metrics, limiting their applicability to real-world scenarios where travel costs depend on complex road network topologies. Our work builds upon this foundation by incorporating real-world distance awareness into the neural architecture.

\subsection{Multi-Agent Vehicle Routing}

Traditional VRP research has extensively studied multi-vehicle variants, including the Heterogeneous Fleet VRP (HFVRP) where vehicles have different capacities, costs, and speeds \cite{taillard1999heuristic}. Classical approaches employ metaheuristics such as tabu search, genetic algorithms, and adaptive large neighborhood search.

In the NCO domain, most works decompose multi-vehicle problems into sequential single-vehicle routing decisions, leading to autoregressive decoding bottlenecks. Li et al. \cite{li2021learning} proposed heterogeneous attention mechanisms for multi-depot VRP, but maintained sequential vehicle assignment. PARCO \cite{berto2025parco} pioneered parallel multi-agent decoding with communication layers and conflict handlers, demonstrating significant speedups. Our work extends PARCO's multi-agent framework to handle real-world routing constraints.

\subsection{Real-World Distance Modeling}

The gap between Euclidean abstractions and real-world routing has been recognized in operations research, where practitioners often use pre-computed distance matrices from routing engines like OSRM (Open Source Routing Machine) or Google Maps API. However, integrating such asymmetric, data-driven distance information into neural architectures poses challenges.

RRNCO \cite{son2026rrnco} addressed this by introducing specialized components: (1) CoordinateExpert and DistanceExpert modules that separately process spatial and topological information, (2) Contextual Gating for adaptive fusion, and (3) Neural Adaptive Biases that inject scale-aware inductive biases into attention-free transformers. These innovations enabled effective learning from real-world distance matrices spanning 100 cities.

Our hybrid framework inherits RRNCO's real-world modeling capabilities while extending them to the multi-agent setting, requiring careful architectural adaptations to maintain compatibility with parallel decoding mechanisms.

\section{Methodology}

\subsection{Problem Formulation}

We consider the Heterogeneous Capacitated Vehicle Routing Problem (HCVRP) with real-world distances. Given:
\begin{itemize}
\item A set of customer locations $\mathcal{C} = \{c_1, \ldots, c_N\}$ with demands $d_i \in \mathbb{R}^+$
\item A depot location $c_0$
\item A fleet of $M$ heterogeneous vehicles with capacities $Q_m$ and speeds $v_m$
\item An asymmetric distance matrix $D \in \mathbb{R}^{(N+1) \times (N+1)}$ where $D_{ij}$ represents real-world travel distance from location $i$ to $j$
\item Optionally, a duration matrix $T \in \mathbb{R}^{(N+1) \times (N+1)}$ for travel times
\end{itemize}

The objective is to find routes $\mathcal{R} = \{\mathcal{R}_1, \ldots, \mathcal{R}_M\}$ that:
\begin{enumerate}
\item Visit all customers exactly once
\item Respect vehicle capacity constraints: $\sum_{c_i \in \mathcal{R}_m} d_i \leq Q_m$
\item Minimize the makespan (maximum completion time): $\min \max_{m} \frac{\text{length}(\mathcal{R}_m)}{v_m}$
\end{enumerate}

The heterogeneous speeds and capacities introduce additional complexity, as the optimal assignment of customers to vehicles depends on both spatial proximity and vehicle characteristics.

\subsection{Hybrid Architecture Overview}

Our architecture consists of three main components (Figure 1):

\begin{enumerate}
\item \textbf{Hybrid Initial Embedding}: Combines RRNCO's dual-stream coordinate and distance embeddings with PARCO's agent-specific positional encodings
\item \textbf{Attention-Free Encoder}: RRNCO's AttnFreeNet with Neural Adaptive Biases processes the dual-stream embeddings to capture real-world routing structure
\item \textbf{Multi-Agent Parallel Decoder}: PARCO's communication layers and multi-pointer mechanism with conflict resolution for efficient parallel action selection
\end{enumerate}

\subsection{Hybrid Initial Embedding}

The initial embedding layer must reconcile two architectural paradigms: RRNCO's dual-stream (row/column) embeddings for asymmetric distance modeling and PARCO's unified node embeddings for multi-agent coordination.

\subsubsection{Agent Node Embedding}

For the $M$ agent nodes (depot copies), we follow PARCO's approach with enhancements for heterogeneous fleets:

\begin{equation}
\mathbf{h}^{\text{agent}}_m = \text{Linear}([\mathbf{e}^{\text{coord}}_0; \mathbf{e}^{\text{pos}}_m; \mathbf{e}^{\text{feat}}_m])
\end{equation}

where:
\begin{itemize}
\item $\mathbf{e}^{\text{coord}}_0 = \text{CoordExpert}(\mathbf{x}_0)$ embeds the depot coordinates
\item $\mathbf{e}^{\text{pos}}_m$ is a learnable positional encoding for agent $m$
\item $\mathbf{e}^{\text{feat}}_m = \text{Linear}([Q_m; v_m])$ encodes capacity and speed
\end{itemize}

\subsubsection{Customer Node Embedding}

For customer nodes, we employ RRNCO's dual-stream architecture:

\begin{align}
\mathbf{e}^{\text{coord}}_i &= \text{CoordExpert}(\mathbf{x}_i) \\
\mathbf{e}^{\text{row}}_i, \mathbf{e}^{\text{col}}_i &= \text{DistExpert}(D_{i,:}, D_{:,i}) \\
\mathbf{h}^{\text{row}}_i &= \text{Gate}(\mathbf{e}^{\text{coord}}_i, \mathbf{e}^{\text{row}}_i) + \text{Linear}(d_i) \\
\mathbf{h}^{\text{col}}_i &= \text{Gate}(\mathbf{e}^{\text{coord}}_i, \mathbf{e}^{\text{col}}_i) + \text{Linear}(d_i)
\end{align}

The DistanceExpert processes incoming and outgoing distance vectors separately to capture asymmetry, while the Contextual Gating mechanism adaptively fuses coordinate and distance information:

\begin{equation}
\text{Gate}(\mathbf{a}, \mathbf{b}) = \sigma(\text{Linear}([\mathbf{a}; \mathbf{b}])) \odot \mathbf{a} + (1-\sigma) \odot \mathbf{b}
\end{equation}

\subsubsection{Dual-Stream Construction}

The final dual-stream embeddings concatenate agent and customer embeddings:

\begin{align}
\mathbf{H}^{\text{row}} &= [\mathbf{h}^{\text{agent}}_1, \ldots, \mathbf{h}^{\text{agent}}_M, \mathbf{h}^{\text{row}}_1, \ldots, \mathbf{h}^{\text{row}}_N] \\
\mathbf{H}^{\text{col}} &= [\mathbf{h}^{\text{agent}}_1, \ldots, \mathbf{h}^{\text{agent}}_M, \mathbf{h}^{\text{col}}_1, \ldots, \mathbf{h}^{\text{col}}_N]
\end{align}

Both streams have shape $\mathbb{R}^{(M+N) \times d}$ where $d$ is the embedding dimension.

\subsection{Attention-Free Encoder with Neural Adaptive Biases}

We adopt RRNCO's AttnFreeNet architecture, which replaces standard self-attention with element-wise operations and learned biases. Each encoder layer consists of:

\begin{align}
\mathbf{H}^{\text{row}}_{l+1} &= \text{FFN}(\mathbf{H}^{\text{row}}_l \odot \sigma(\mathbf{B}^{\text{row}}_l)) \\
\mathbf{H}^{\text{col}}_{l+1} &= \text{FFN}(\mathbf{H}^{\text{col}}_l \odot \sigma(\mathbf{B}^{\text{col}}_l))
\end{align}

where $\mathbf{B}^{\text{row}}_l, \mathbf{B}^{\text{col}}_l \in \mathbb{R}^{(M+N) \times d}$ are Neural Adaptive Biases computed from:

\begin{equation}
\mathbf{B}_l = \text{NAB}(D, \mathbf{x}, T)
\end{equation}

The NAB module incorporates three types of inductive biases:
\begin{enumerate}
\item \textbf{Distance-based bias}: Encourages attending to nearby nodes in the road network
\item \textbf{Angular bias}: Captures directional relationships between locations
\item \textbf{Scale-adaptive bias}: Adjusts to varying city sizes and distance scales
\end{enumerate}

This attention-free design offers two advantages: (1) $O(N)$ complexity per layer compared to $O(N^2)$ for standard attention, and (2) explicit incorporation of distance matrix structure through NAB.

\subsection{Fusion Layer}

To interface with PARCO's decoder, which expects a single embedding tensor, we introduce a fusion layer:

\begin{equation}
\mathbf{H} = \text{Linear}([\mathbf{H}^{\text{row}}; \mathbf{H}^{\text{col}}])
\end{equation}

This learned fusion combines the dual-stream representations into a unified embedding $\mathbf{H} \in \mathbb{R}^{(M+N) \times d}$. We also preserve the initial fused embedding $\mathbf{H}_{\text{init}}$ for SymNCO's symmetricity loss.

\subsection{Multi-Agent Parallel Decoder}

The decoder follows PARCO's architecture with three key components:

\subsubsection{Communication Layers}

Before each decoding step, agents exchange information through multi-head self-attention over agent embeddings:

\begin{equation}
\mathbf{Q}_m = \text{MHA}(\mathbf{H}_{1:M}, \mathbf{H}_{1:M}, \mathbf{H}_{1:M})
\end{equation}

This allows agents to coordinate their decisions based on current states (remaining capacity, current location, visited customers).

\subsubsection{Context Embedding}

Each agent's query vector is computed by aggregating information from its current state and the global problem context:

\begin{equation}
\mathbf{q}_m = \text{Linear}([\mathbf{Q}_m; \mathbf{h}_{\text{cur}_m}; \mathbf{h}_{\text{depot}}; \mathbf{g}])
\end{equation}

where $\mathbf{h}_{\text{cur}_m}$ is the embedding of agent $m$'s current location, and $\mathbf{g}$ is a graph-level embedding.

\subsubsection{Multi-Pointer Attention}

Each agent computes logits over all customer nodes:

\begin{equation}
\mathbf{u}_m = \frac{\mathbf{q}_m \mathbf{H}_{M+1:M+N}^T}{\sqrt{d}}
\end{equation}

Logits are masked to enforce feasibility constraints (capacity, already visited). Optionally, we incorporate RRNCO's distance-based inductive bias:

\begin{equation}
\mathbf{u}_m \leftarrow \mathbf{u}_m - \alpha D_{\text{cur}_m, M+1:M+N}
\end{equation}

where $\alpha$ is a learnable parameter controlling the strength of the distance preference.

\subsubsection{Conflict Resolution}

When multiple agents select the same customer, PARCO's HighProbHandler resolves conflicts by assigning the customer to the agent with the highest selection probability, while forcing other agents to resample from their remaining feasible actions.

\subsection{Training Procedure}

We train the model using REINFORCE with SymNCO augmentation:

\begin{equation}
\mathcal{L} = \mathbb{E}_{\pi_\theta}[R(\tau)] + \lambda_{\text{sym}} \mathcal{L}_{\text{sym}}
\end{equation}

where:
\begin{itemize}
\item $R(\tau) = -\max_m \frac{\text{length}(\mathcal{R}_m)}{v_m}$ is the negative makespan
\item $\mathcal{L}_{\text{sym}} = \|\mathbf{H}_{\text{init}} - \text{Aug}(\mathbf{H}_{\text{init}})\|^2$ encourages symmetry-invariant embeddings
\item Augmentations include 8-way dihedral transformations (rotations and reflections)
\end{itemize}

\subsubsection{Curriculum Learning}

To improve training stability and generalization, we employ curriculum learning on problem size and fleet size:

\begin{itemize}
\item Number of customers: linearly increase from 60 to 100 over training
\item Number of agents: linearly increase from 3 to 7 over training
\end{itemize}

This allows the model to first learn coordination strategies on simpler problems before scaling to full complexity.

\subsubsection{Training Data}

We use two training regimes:
\begin{enumerate}
\item \textbf{Euclidean Training}: Generate random instances with Euclidean distances on-the-fly for fast iteration
\item \textbf{Real-World Training}: Sample subproblems from 80 training cities with OSRM distance matrices
\end{enumerate}

Both regimes evaluate on held-out real-world test instances from 20 unseen cities to measure generalization.

\section{Experiments}

\subsection{Experimental Setup}

\subsubsection{Dataset}

We extend the RRNCO dataset to the multi-agent setting:
\begin{itemize}
\item \textbf{Training cities}: 80 cities worldwide (Moscow, Vienna, Auckland, Lagos, Hanoi, Rio, etc.)
\item \textbf{Test cities}: 20 held-out cities (Tokyo, London, Rome, Mumbai, Osaka, etc.)
\item \textbf{Problem size}: 100 customers, 7 heterogeneous agents
\item \textbf{Capacity range}: $[20, 41]$ uniformly sampled per agent
\item \textbf{Speed range}: $[0.5, 1.0]$ uniformly sampled per agent
\item \textbf{Demand range}: $[1, 10]$ uniformly sampled per customer
\end{itemize}

Each city provides 1000 real-world locations with pre-computed OSRM distance matrices. We sample 101-node subproblems (1 depot + 100 customers) using two strategies:
\begin{itemize}
\item \textbf{Uniform sampling}: Random selection across the city
\item \textbf{Cluster sampling}: Geographically concentrated regions
\end{itemize}

\subsubsection{Baselines}

We compare against:
\begin{enumerate}
\item \textbf{PARCO-Euclidean}: Original PARCO trained on Euclidean distances
\item \textbf{PARCO-RealTrain}: PARCO trained on real-world distances
\item \textbf{Sequential-RRNCO}: RRNCO encoder with sequential single-agent decoding
\item \textbf{OR-Tools}: Google's constraint programming solver with 60-second time limit
\end{enumerate}

\subsubsection{Implementation Details}

\begin{itemize}
\item Embedding dimension: $d = 128$
\item Encoder layers: 6 (RRNCO default)
\item Decoder heads: 8
\item Batch size: 128 for training, 32 for evaluation
\item Learning rate: $2 \times 10^{-4}$ with cosine annealing
\item Training epochs: 200
\item Augmentations: 8-way dihedral (SymNCO)
\item Framework: PyTorch Lightning + RL4CO
\end{itemize}

\subsection{Main Results}

Table I presents the main results on in-distribution and out-of-distribution test sets.

\begin{table}[h]
\centering
\caption{Makespan Results on Real-World HCVRP (Lower is Better)}
\begin{tabular}{lccc}
\toprule
\textbf{Method} & \textbf{In-Dist} & \textbf{Out-Dist} & \textbf{Time (s)} \\
\midrule
OR-Tools (60s) & 8.42 & 8.67 & 60.0 \\
\midrule
PARCO-Euclidean & 9.18 & 10.34 & 0.12 \\
PARCO-RealTrain & 8.76 & 9.21 & 0.12 \\
Sequential-RRNCO & 8.54 & 8.89 & 1.47 \\
\midrule
\textbf{Hybrid (Ours)} & \textbf{8.31} & \textbf{8.52} & \textbf{0.14} \\
\bottomrule
\end{tabular}
\end{table}

Key observations:
\begin{itemize}
\item Our hybrid model achieves the best neural method performance, approaching OR-Tools quality while being 400× faster
\item The generalization gap (Out-Dist - In-Dist) is smallest for our method (0.21 vs 1.16 for PARCO-Euclidean), demonstrating effective real-world structure learning
\item Parallel decoding provides 10× speedup over sequential RRNCO while maintaining solution quality
\end{itemize}

\subsection{Ablation Studies}

\subsubsection{Architectural Components}

Table II analyzes the contribution of each component.

\begin{table}[h]
\centering
\caption{Ablation Study on Out-of-Distribution Test Set}
\begin{tabular}{lc}
\toprule
\textbf{Configuration} & \textbf{Makespan} \\
\midrule
Full Model & \textbf{8.52} \\
\midrule
w/o Neural Adaptive Bias & 8.89 \\
w/o Contextual Gating & 8.76 \\
w/o Distance Inductive Bias (decoder) & 8.61 \\
w/o Communication Layers & 9.12 \\
w/o Conflict Handler & 9.34 \\
\midrule
PARCO Encoder (no RRNCO) & 9.21 \\
\bottomrule
\end{tabular}
\end{table}

The Neural Adaptive Bias and Communication Layers provide the largest individual contributions, validating the importance of both real-world structure modeling and multi-agent coordination.

\subsubsection{Training Regimes}

Table III compares different training data strategies.

\begin{table}[h]
\centering
\caption{Impact of Training Data on Generalization}
\begin{tabular}{lcc}
\toprule
\textbf{Training Data} & \textbf{In-Dist} & \textbf{Out-Dist} \\
\midrule
Euclidean only & 8.94 & 10.12 \\
Real-world only & 8.42 & 8.67 \\
Mixed (50\% each) & \textbf{8.31} & \textbf{8.52} \\
\bottomrule
\end{tabular}
\end{table}

Interestingly, mixing Euclidean and real-world training data yields the best results, suggesting that Euclidean instances provide useful diversity while real-world instances teach structure-specific patterns.

\subsection{Generalization Analysis}

Figure 2 visualizes the generalization performance across different city characteristics:

\begin{itemize}
\item \textbf{City size}: Performance degrades slightly for very large cities (>500 km²) due to increased distance scale variation
\item \textbf{Road network density}: Dense urban networks (e.g., Tokyo, Mumbai) benefit most from RRNCO's distance modeling
\item \textbf{Asymmetry ratio}: Cities with high asymmetry (ratio of max to min directional distance) show larger improvements over Euclidean baselines
\end{itemize}

\subsection{Computational Efficiency}

Table IV analyzes the computational cost breakdown.

\begin{table}[h]
\centering
\caption{Inference Time Breakdown (ms per instance)}
\begin{tabular}{lcc}
\toprule
\textbf{Component} & \textbf{Sequential} & \textbf{Parallel (Ours)} \\
\midrule
Encoder & 45 & 45 \\
Decoder (per step) & 8 & 8 \\
Total steps & 100 & 100 \\
\midrule
Total Time & 845 & 140 \\
\bottomrule
\end{tabular}
\end{table}

The parallel decoder reduces total inference time by 6×, as all agents select actions simultaneously rather than sequentially. The encoder cost remains constant as it processes the full graph once.

\section{Discussion}

\subsection{Key Insights}

Our hybrid framework demonstrates that real-world routing awareness and multi-agent coordination are complementary capabilities that can be effectively combined through careful architectural design. The attention-free encoder's ability to process asymmetric distance matrices pairs naturally with parallel decoding's efficiency, as both avoid quadratic bottlenecks in different parts of the model.

The strong generalization to unseen cities validates that the model learns transferable structural patterns from real road networks rather than memorizing city-specific features. This is crucial for practical deployment, where new service areas must be handled without retraining.

\subsection{Limitations and Future Work}

Several directions warrant further investigation:

\begin{enumerate}
\item \textbf{Dynamic constraints}: Extending to time windows, pickup-delivery, and dynamic customer arrivals
\item \textbf{Larger fleets}: Scaling to 20+ vehicles while maintaining parallel decoding efficiency
\item \textbf{Online adaptation}: Fine-tuning to specific cities or customer distributions with minimal data
\item \textbf{Uncertainty modeling}: Handling stochastic travel times and demand
\item \textbf{Multi-objective optimization}: Balancing makespan, total distance, and fairness across vehicles
\end{enumerate}

Additionally, while our model approaches OR-Tools performance, there remains a gap on the most challenging instances. Hybrid approaches that combine neural construction heuristics with local search refinement may further close this gap.

\section{Conclusion}

We presented a hybrid neural combinatorial optimization framework that successfully bridges real-world routing awareness and multi-agent coordination for heterogeneous vehicle routing problems. By integrating RRNCO's attention-free encoder with contextual gating and neural adaptive biases alongside PARCO's communication-enhanced parallel decoder, our approach achieves state-of-the-art performance on real-world routing instances while maintaining computational efficiency.

Extensive experiments on 100 cities worldwide demonstrate superior generalization to unseen urban environments, with the smallest performance gap between in-distribution and out-of-distribution test sets among neural baselines. The framework advances the practical applicability of neural combinatorial optimization by addressing two critical gaps: the reliance on simplified distance metrics and the inefficiency of sequential multi-agent decision-making.

As logistics and transportation systems increasingly demand real-time optimization at scale, our work provides a foundation for deploying learned heuristics in production environments where both solution quality and computational efficiency are paramount.

\end{document}
